\emph{The hypothesis failed on the passive AFO, for a reason the hypothesis did not anticipate.} A passive AFO is expected to cost push-off because it resists plantarflexion. With its neutral angle at 0 and a stiffness of \SI{160}{N.m/rad}, the spring chosen by the rule instead stores energy in the stance dorsiflexion and returns it in early push-off, so the total ankle power rises while the muscles' share falls by 9\%. Whether a passive AFO costs push-off depends on its neutral angle and on how far the ankle dorsiflexes before push-off, not only on its stiffness. A clinical posterior leaf-spring AFO is set near neutral too, but the model's stance dorsiflexion of up to \SI{10}{\degree} under a stiff spring is more than an orthosis with soft-tissue compliance would see. The result is a statement about this metric and this spring, not evidence that passive AFOs help push-off.

\emph{Toe clearance was the wrong primary metric for an adapted model.} The adapted model already cleared its toe 9 to 20~mm \emph{higher} than healthy by lifting the hip, so part 1 could only fail if a device made clearance worse. It did not discriminate between the devices: all three re-optimized conditions have the same mean clearance. What the active AFO changed was how the clearance was produced: with the foot held up, the re-optimized controller gave up most of the steppage, and the cost of transport fell by 5\%, about two thirds of the way back to healthy. Steppage or the cost of transport would have been the more telling primary metric, and the toe clearance threshold should have been two-sided. Both choices are recorded here as lessons, not applied after the fact.

\emph{Adaptation dominates.} No device setting could be worn by the frozen controllers, and even a \SI{5}{N.m/rad} spring made them fall. The comparison between devices is therefore only meaningful after re-optimization, which assumes the wearer re-tunes the whole gait to the device. People adapt partly and over time; the truth for a first session lies between a model that falls and a model that has fully re-optimized. Device studies on optimized reflex models should report both, as here, and should not read the frozen results as a prediction of a person's first steps.

\emph{The sensing and actuation chain was not the bottleneck.} The shank-gyroscope detector found 99\% of the events on drop-foot gait with heel-strike errors of about 24~ms, and transferred to real gyroscopes without retuning. Doubling its noise changed nothing. The SEA tracked the requirement to 1.5\% with no saturation, and the simulated device used the slower fitted response without the model feedforward. The fragility of the re-optimized gaits to timing (one seed falls at any IMU or actuator delay other than the nominal one) comes from the optimized gait, not from the device: the active AFO itself still lifted the foot at 100~ms of delay. For the drop-foot deficit the series spring saved almost no energy; its value in such a device is torque sensing and impact tolerance.
